\section{Alarm Mechanism and Gap Decomposition}
\label{sec:gap_decomposition}

\subsection{The alarm}

\begin{definition}[B-to-C Alarm Under Revealed Sacrifice]
\label{def:alarm}
The B-to-C alarm fires when the lower-bound B-to-C ratio drops
below a threshold:
\[
  \mathrm{ALARM}(T) \;=\; \betaL(G, T) \;<\; \betaalarm,
\]
where $\betaalarm \in (0, 1)$ is configurable.  The alarm detects
proxy failure: when $\betaL$ is low, $\volL$ is high but the
sacrifice channel shows little evidence of the corresponding
$\volR$.
\end{definition}

\begin{remark}[Alarm conservatism]
Because $\betaL \leq \betatrue$, the alarm may fail to fire when
$\betatrue < \betaalarm$ but the sacrifice channel has
insufficient coverage.  The alarm has no false positives (if
$\betaL$ is low, $\volR$ is genuinely not well-evidenced) but
may have false negatives ($\volR$ may be high but invisible to
the sacrifice channel).  This is the cost of privacy-minimal
observation.
\end{remark}

\subsection{Diagnostic decomposition}

\begin{definition}[Gap Decomposition Under Revealed Sacrifice]
\label{def:gap_decomp}
The B-to-C gap ($1 - \betaL$) decomposes into four \emph{disjoint}
sources.  Each capability in $P$ is assigned to exactly one
category by priority ordering (restricted $>$ covered $>$ dormant
$>$ residual), and each category's $\delta$ term is its uncovered
$\volL$ share:
\[
  1 - \betaL \;=\; \delta_{\mathrm{restricted}}
  + \delta_{\mathrm{partial}} + \delta_{\mathrm{dormant}}
  + \delta_{\mathrm{residual}},
\]
where the capability partition is:
\begin{itemize}
\item \textbf{Restricted} ($\Xi$): $d$ is currently restricted
  under the preemptive-restriction criterion
  of~\citep[Proposition~1]{lasser2026horizon}.
  Classified first regardless of sacrifice evidence.
\item \textbf{Covered} ($\CovS$): $d \notin \Xi$, and $d$
  appears in at least one acquired bundle $Y_n$ for some
  theorem-admissible sacrifice event in $[T - H, T]$.  A
  theorem-admissible event satisfies S0, S1, and S2.
\item \textbf{Dormant} ($\DorS$): $d \notin \Xi$,
  $d \notin \ResS$, and $d$ does not appear in any acquired
  bundle of any theorem-admissible event in $[T - H, T]$.
  These are tradable-in-principle capabilities with zero
  want-pursuit sacrifice evidence---dormant \emph{wants}.
\item \textbf{Residual} ($\ResS$): $d \notin \Xi$, and $d$ is
  structurally outside the sacrifice channel
  (Definition~\ref{def:residual_class}).
\end{itemize}

The $\delta$ terms are:
\begin{align*}
  \delta_{\mathrm{restricted}} &= [\volL(\Xi) -
    \volRlower(\Xi)] \;/\; \volL(G), \\
  \delta_{\mathrm{partial}} &= [\volL(\CovS) -
    \volRlower(\CovS)] \;/\; \volL(G), \\
  \delta_{\mathrm{dormant}} &= \volL(\DorS) \;/\; \volL(G), \\
  \delta_{\mathrm{residual}} &= \volL(\ResS) \;/\; \volL(G).
\end{align*}
\end{definition}

\begin{proposition}[Computability of Gap Decomposition]
\label{prop:gap_computable}
The four-term gap decomposition is computable from:
the aggregate sacrifice data (events in $[T - H, T]$),
the restriction set $\Xi$,
the residual class $\ResS$
(Definition~\ref{def:residual_class}), and
the capability census (enumeration of $P$).
Classification is $O(|P|)$ given the sacrifice database and the
$\ResS$ classification.  No structural counterfactual computation
is required (a direct-measurement approach would require
$O(|P|^2)$ counterfactual queries).
\end{proposition}

\begin{remark}[Part~A versus Part~B granularity]
Under Part~A alone (category-level aggregates, S0--S2), the
decomposition operates at the category level: each partition
component $C_j$
(Definition~\ref{def:category_partition}) is classified as
restricted, covered, dormant, or residual.  Under committed data
(Section~\ref{sec:commitment}), S2 is not directly verifiable;
the committed-mode decomposition therefore \emph{assumes} S2
holds for committed events.  The per-capability decomposition is
the refinement available when Part~B assumptions hold.
\end{remark}

\subsection{Need-sufficiency diagnostic (companion metric)}
\label{sec:need_diagnostic}

The gap decomposition above tracks want-pursuit: how much of the
collective's above-sufficiency capability-space is evidenced by
voluntary sacrifice.  The need-sufficiency diagnostic is the
parallel metric for below-sufficiency: how efficiently does the
collective satisfy needs, and is the satisfaction downstream-safe?

\begin{definition}[Need-Access Cost]
\label{def:nac}
The need-access cost aggregate for trade window $[T - H, T]$ is:
\[
  \NAC(T) \;=\; \sum_{n \in F_{S1}}
  \Delta\volL(X_n),
\]
where $F_{S1}$ is the set of events failing S1 due to
below-sufficiency (Proposition~\ref{prop:s1_sufficiency}).
$\NAC$ measures the total benchmarked capability surrendered for
need-satisfaction---purchasing power and time spent reaching
baseline functionality rather than pursuing wants.

$\NAC$ does \emph{not} enter
Theorem~\ref{thm:aggregate_bound}'s aggregate lower bound or
$\betaL$.  It is a separate diagnostic that tracks the cost side
of the agent's capability budget.
\end{definition}

\begin{remark}[$\NAC$ is an upper bound on need-access cost]
Because the whole-event filter sends the entire
$\Delta\volL(X_n)$ of mixed need/want events into $\NAC$,
$\NAC$ overestimates the true need-access cost by including any
attached want-pursuit spend.  $\NAC$ should be read as an
\emph{upper bound} on need-access cost, not a precise
measurement.
\end{remark}

\begin{definition}[Need-Satisfaction Efficiency]
\label{def:nse}
For each need $N$ identified in the poset
(Definition~\ref{def:need_bundle}) and each agent~$i$, the
per-agent need-satisfaction efficiency is a triple:
\[
  \NSE(i, N) \;=\; (\mathrm{sufficiency\_status},\;
  \mathrm{downstream\_safety},\; \mathrm{access\_cost}),
\]
where:
\begin{itemize}
\item $\mathrm{sufficiency\_status} \in
  \{\text{below}, \text{met}, \text{exceeded}\}$;
\item $\mathrm{downstream\_safety} \in
  \{\text{safe}, \text{degrading}, \text{unknown}\}$;
\item $\mathrm{access\_cost} = \sum \Delta\volL(X_n)$ for
  below-sufficiency events by agent~$i$ targeting need~$N$.
\end{itemize}
The collective-level aggregate $\NSE(T)$ reports the distribution
of (agent, need) pairs across sufficiency statuses and flags
degrading pathways.
\end{definition}

\begin{remark}[Observability under commitment]
When sacrifice events are committed
(Section~\ref{sec:commitment}), per-agent $\NSE(i, N)$ is not
recoverable.  The committed-mode fallback is a population-level
proxy: $\NSE_{\mathrm{proxy}}(T)$ reports $\NAC(T)$ by need
category and counts of below-sufficiency events.  Full $\NSE(T)$
requires agents to voluntarily attest pathways or the framework to
operate in non-committed mode.
\end{remark}

\begin{proposition}[Combined Diagnostic]
\label{prop:combined_diagnostic}
The framework reports the triple $(\volRlower, \NAC, \NSE)$ as a
joint diagnostic:
\begin{enumerate}
\item[(a)] \textbf{High $\volRlower$ + low $\NAC$:} Healthy.
  The collective exercises rich capabilities above sufficiency
  and needs are cheap to meet.
\item[(b)] \textbf{Low $\volRlower$ + high $\NAC$:} Proxy
  failure on needs.  $\volL$ may be high, but agents are
  spending most of their benchmarked capability budget on basic
  need-satisfaction.  The sacrifice channel correctly excludes
  these events (S1 fails), so the low $\volRlower$ accurately
  reflects low above-sufficiency exercise.  $\NAC$ surfaces
  \emph{why} $\volRlower$ is low.
\item[(c)] \textbf{Low $\volRlower$ + low $\NAC$:} Dormancy,
  residual, or theorem-inadmissible activity.  The gap
  decomposition (Definition~\ref{def:gap_decomp}) distinguishes
  dormant from residual.
\item[(d)] \textbf{High $\volRlower$ + high $\NAC$:} Mixed
  regime.  $\NSE$ distinguishes ``expensive but safe'' from
  ``expensive and degrading.''
\end{enumerate}
\end{proposition}

\begin{remark}[Parallel structure, not nesting]
The need-sufficiency diagnostic is parallel to the gap
decomposition, not nested within it, but operates on different
formal objects.  The gap decomposition is a \emph{capability
partition}; $\NAC$ is an \emph{event aggregate}; $\NSE$ is a
\emph{per-need state assessment}.  The two diagnostics cover
complementary portions of the sacrifice event stream:
S1-satisfying events (that also satisfy S0 and S2) enter the gap
decomposition via Theorem~\ref{thm:aggregate_bound}; events
failing S1 due to below-sufficiency enter $\NAC$.  Events failing
S1 for non-need reasons (coercion, duress) are excluded from
both diagnostics and constitute an undiagnosed event-stream
residual.
\end{remark}
